% !TeX program = xelatex
\documentclass{jos}
\fancyhead[LE,RO]{\joskaiti HyMEM 阅读报告}
\fancyhead[LO,RE]{\joskaiti 第8队\quad lanshuheng9-oss}
\fancyfoot[C]{\small\thepage}
\setlength{\headheight}{16pt}
\hypersetup{pdftitle={HyMEM 阅读报告：录音配套修订},pdfauthor={lanshuheng9-oss}}
\begin{document}
\raggedbottom
\begin{center}
{\Large\bfseries HyMEM 阅读报告}\par
第8队\quad lanshuheng9-oss\quad 2026年10月27日
\end{center}
\section{HyMEM：电脑操作智能体的混合自演化记忆}
\label{sec:hymem}

HyMEM\upcite{zhu-etal-2026-hybrid}面向需要观察屏幕截图并连续操作网页的电脑操作智能体，通过融合离散策略与连续轨迹表示构建混合记忆网络，并配合跨任务的长期记忆演化与任务内的工作记忆刷新，动态适应交互阶段的变化。

\subsection{问题背景与混合表示}

网页操作智能体以任务描述、当前截图和历史成功轨迹为输入，逐步生成下一步可执行动作。在此类多模态任务中，历史经验兼含高层策略与底层界面交互细节：纯文本总结容易丢失按钮位置、页面布局等视觉线索，单纯连续编码又缺乏清晰的策略组织与可解释索引；同时，随着任务从商品搜索推进至结账，子目标已发生转移，初始检索的经验不再适用，但全局预算等约束仍需保留。

论文在记忆图 $G=(V,E)$ 中以三元组表示单条历史轨迹：
\begin{equation}
v_i=(c_i,A_i,m_i),
\end{equation}
其中 $c_i$ 为高层策略，$A_i$ 为动作、界面元素或领域标签构成的属性集合，$m_i$ 为连续轨迹编码。共享相同属性的节点通过无向边相连，利用离散策略与属性构建语义索引网络；同时借鉴CoMEM的连续编码器将每条多模态轨迹压缩为8个嵌入向量 $m_i$，保留难以用文字概括的视觉与动作细节。检索时，系统将归一化后的任务文本CLIP向量与起始截图CLIP向量拼接，通过FAISS寻找相似经验。

\noindent\begin{minipage}{\linewidth}
\setlength{\parindent}{2em}
\subsection{全局记忆演化与结构化检索}

更新长期记忆图时，系统先为新轨迹检索相似候选经验，再由视觉语言模型评估其信息增量，并执行三类操作：

（1）ADD：出现新策略或新属性时，新增对应节点与连接。

（2）MERGE：已有策略获得界面变体或替代路径等互补证据时，将新证据并入对应策略及属性关联。

（3）REPLACE：在同等条件下，新轨迹以更少步骤、更高成功率或更低风险实现同一策略时，用其替换过时证据并刷新关联。

上述操作在长期图中持续积累可迁移经验并抑制冗余。

\end{minipage}\par

推理检索时，系统先依据任务与截图召回语义相似的种子轨迹，再沿共享属性边扩展邻居并按相似度排序。主实验每次取回10条轨迹，由5条语义相似种子及其1跳邻居中选出的5条关联轨迹组成。

\subsection{混合工作记忆及局部刷新}

召回的经验形成两个互补视图：策略与属性汇总为指导指令（Guidance Instructions）写入系统提示词，辅助宏观规划；连续轨迹嵌入 $m_i$ 则直接接入模型输入，提供细粒度视觉与动作参照，二者共同构成当前任务的工作记忆。

每步动作后，视觉语言模型对比前后页面观测与当前指导，判断是否发生阶段转换。若任务由搜索转入结账，系统在保留全局目标与约束的前提下剔除过时的阶段指导，再按新页面状态重新检索，合并新指导并刷新连续嵌入，从而在任务内部实现工作上下文的动态更新。

\subsection{实验设置与结果}

论文在WebVoyager、Multimodal-Mind2Web和MMInA三个多模态网页智能体基准上评测。前两者根据任务描述与截图轨迹由模型判分，MMInA比对最终回答与标准答案，主指标均为各领域平均任务成功率（Overall）。记忆库由GUIAct、Mind2Web训练集及智能体自身生成的2,883条成功轨迹构成，主实验图含1,858个节点和逾100万条边。执行采用ReAct接口，配合SOM界面标记及UI-INS-7B定位回退；训练仅更新记忆编码器的Q-Former与LoRA层（约占总参数的1.2\%，训练1个epoch）。

\begin{table}[!htbp]
\centering
\caption{不同记忆表示下的Overall任务成功率（\%）}
\label{tab:hymem-results}
\begin{tabular}{lrrrr}
\toprule
基础模型 & 无记忆 & 离散记忆 & 连续记忆 & HyMEM \\
\midrule
Qwen2.5-VL-7B & 12.5 & 17.0 & 21.7 & 35.0 \\
Qwen3-VL-8B & 21.9 & 21.3 & 23.3 & 33.8 \\
UI-TARS-1.5-7B & 17.1 & 15.4 & 28.7 & 31.2 \\
\bottomrule
\end{tabular}
\end{table}

表\ref{tab:hymem-results}显示，HyMEM在三种模型上均取得最高总体成功率，其中Qwen2.5-VL-7B从无记忆时的12.5\%升至35.0\%（提升22.5个百分点），显著优于单独的离散或连续记忆，验证了高层策略与轨迹细节的互补性。后两种模型在仅使用离散记忆时略低于无记忆基线，结合连续轨迹后则均获明显提升。在Amazon与Google Maps的消融分析中，相较于静态记忆，全局演化使成功率分别由48.8\%升至61.0\%、43.9\%升至51.2\%，局部刷新亦分别达到56.1\%和48.8\%，证实了两层动态更新机制的作用。

\subsection{方法局限与课设启发}

HyMEM通过混合表示与两层演化提升了网页操作表现，但在记忆构建、编码器训练、图扩展及每步阶段判定上增加了计算开销；属性标签或阶段判断失误也可能引起错误传播。对纯文本对话记忆而言，可重点借鉴其“语义种子＋结构邻居”的扩展检索方式，以及区分“跨会话长期记忆演化”与“当前会话工作上下文刷新”的双层维护思路。

% 原文定位：PDF第1--3页，表1；PDF第3--5页，\S3.1--3.3、图1；PDF第3页，式(1)；PDF第4页，图1与Memory Update；PDF第4--5页，Structured Retrieval及Memory Settings；PDF第5页，On-the-fly Working Memory Refresh；PDF第5--7页，\S4.1--4.2与表2；PDF第5页，\S4.1；PDF第7页，\S4.3.1、图2。


\subsection{结果解释与方法取舍}

在表\ref{tab:hymem-results}中，相较于仅使用连续记忆，HyMEM在Qwen2.5-VL-7B、Qwen3-VL-8B和UI-TARS-1.5-7B上分别进一步提升13.3、10.5和2.5个百分点。这表明将高层文本策略与底层连续轨迹结合普遍优于单一模态记忆，不过不同底座模型的受益幅度有所差异；其中后两种模型在仅使用离散记忆时表现略低于无记忆基线，说明文本化策略的抽象粒度及其与底座模型的适配程度仍会影响实际效果。

以“商品搜索转入下单结账”的典型流程为例，可以直观理解三层表示与两层更新的协同分工：高层策略记录“先搜索、再按条件筛选、最后结算”的通用流程，属性标签索引搜索、筛选与价格等要素，连续轨迹嵌入则保留具体的按钮位置与页面跳转细节。当页面由商品列表切换至结账界面时，系统的局部刷新机制保留用户的全局预算约束，剔除已完成的搜索阶段指导，并根据结账页面重新检索支付与地址填写经验。

受此启发，在课程设计的纯文本对话记忆实验中，可以在固定回答模型与总检索条数的前提下，对比“单纯召回10条语义相似种子”与“召回5条语义种子＋扩展5条结构关联邻居”两种检索策略，通过同时观测关键证据覆盖率、无关噪声比例、回答准确率及调用开销，定量评估结构化邻居扩展的净收益。
% 补充分析依据：HyMEM PDF第3--5页的方法、PDF第7页表2；差值为对原表的算术整理。
\bibliographystyle{jos}
\bibliography{references}
\end{document}
